\chapter{Tableaux de bord par puce}\label{chap:dashboards}

L'oscilloscope vous montre la forme du son ; le tableau de bord, en dessous, vous dit comment la puce la fabrique.
Nous n'avons pas voulu une vue unique remplie de registres : chaque famille de puces a la sienne, construite autour de
ce qui fait son son. Sur le SID, c'est le filtre et les modulations entre voix ; sur la Super Nintendo, les
enveloppes et l'écho ; sur la PC Engine, la table d'ondes ; sur les puces FM, les opérateurs et la façon dont ils
sont branchés.

Nous relevons l'état de la puce environ toutes les millisecondes de son émulé, et chaque image affiche le dernier
relevé, tel quel. Si une valeur saute, c'est que la puce a sauté.

\section{SID (Commodore 64)}
\capture[0.8]{sidcells-now-playing}{Tableau de bord SID}
Chaque colonne correspond à une voix. De haut en bas, vous lisez :
\begin{description}
  \item[Wave] la forme d'onde (triangle, dent de scie, impulsion, bruit, ou une combinaison) et, pour l'impulsion,
        son rapport cyclique. Quand le rapport cyclique bouge sans arrêt, c'est la « PWM » qui donne leur
        épaisseur aux basses de Rob Hubbard.
  \item[Env] l'enveloppe : le niveau actuel, la phase (attaque, déclin, maintien, relâchement) et les durées ADSR en
        millisecondes. Une attaque à 2~ms et un déclin court, c'est une note percussive ; un maintien haut, une note
        tenue.
  \item[Freq] la note la plus proche et la fréquence en hertz. Un arpège se voit ici : la note change à chaque image.
  \item[Mod] la synchronisation et la modulation en anneau. Nous n'indiquons que la voix source : « Sync V3 »,
        « Ring V2 », « Sync+Ring V1 ». Sur le SID, une voix est toujours modulée par la précédente (V1 par V3, V2 par
        V1, V3 par V2) ; connaître la source suffit.
\end{description}
La dernière ligne montre le volume général, la résonance, le type de filtre et les voix qui passent dedans. Si vous
voyez une voix entrer et sortir du filtre au rythme de la musique, vous tenez l'un des trucs favoris des musiciens
de l'époque.
\capturedouble{sidcells-dual}{sidcells-fullscreen}{Deux puces SID}{plein écran}

\section{Super Nintendo (S-DSP)}
\capturedouble{snes}{snes-inspector}{Vue des 8 voix}{inspecteur d'une voix}
Le S-DSP joue huit voix d'échantillons. Pour chacune, vous voyez le numéro de l'échantillon, la note, l'enveloppe
dessinée (ADSR ou mode GAIN, que beaucoup de drivers utilisent pour leurs fondus maison), les volumes gauche et
droite, et trois indicateurs : écho, bruit, modulation de hauteur.

En bas, les volumes principaux, le volume et le retour de l'écho, et les huit coefficients du filtre FIR qui colore
ses répétitions. Un retour élevé avec un FIR passe-bas, c'est la grande salle réverbérante des jeux Square ; un écho
court et sec, une ambiance de donjon. Touchez une voix pour ouvrir son inspecteur et voir son enveloppe en grand.

\section{PC Engine (HuC6280)}
\capture[0.8]{pce-dashboard}{Vue PC Engine}
Chacune des six voix lit en boucle une table de 32 valeurs de 5 bits : sa forme d'onde. Le rectangle de chaque voix
dessine la table en cours d'utilisation, pas à pas. À côté, une barre de niveau et deux barres de panoramique, puis
la note, la fréquence et le niveau.

Nous décrivons la forme en mots : « Sine », « Triangle », « Saw $\uparrow$ », « Square », « Pulse 25\% »\dots{} et
nous numérotons chaque table distincte. Quand le numéro change en plein morceau, le driver vient de recharger la
voix avec un autre instrument. Le HuC6280 ne sait pas lire une table à l'envers : une scie descendante est une table
que le driver a écrite à l'envers, d'où la mention « Saw $\downarrow$ (reversed) ».

La ligne du bas de chaque voix donne ses registres tels qu'ils ont été écrits : \$802 et \$803 (fréquence),
\$804 (volume et activation), \$805 (panoramique), \$806 (mode DDA), \$807 (bruit). La ligne globale donne \$801
(volume principal), \$808 et \$809 (LFO). Le mode DDA, quand il est actif, montre que la voix joue des échantillons
envoyés un par un par le processeur : c'est ainsi que la PC Engine fait parler ses personnages.

\section{Puces FM}
\capturedouble{fm-megadrive-cards}{fm-megadrive-inspector}{Mega Drive : cartes des voix}{inspecteur d'une voix}
Chaque carte montre une voix : son algorithme, dessiné avec les porteuses en bas et les modulateurs au-dessus, et le
niveau de chaque opérateur. Retenez une règle simple : plus un modulateur est fort, plus le son est riche et dur.
Une basse de Mega Drive qui « claque » a un modulateur dont l'enveloppe chute très vite.

Touchez une voix pour ouvrir l'inspecteur : enveloppes de chaque opérateur dessinées en grand, multiplicateurs de
fréquence, désaccord, rétroaction de l'opérateur 1 et sensibilité au LFO.
\capturedouble{fmview-algorithms}{fm-adlib-cards}{Algorithmes FM}{AdLib}
Pour l'AdLib (OPL2), chaque voix n'a que deux opérateurs et deux branchements possibles : en série (l'un module
l'autre) ou en parallèle (les deux sont entendus). Nous disposons les neuf voix en 3 sur 3.

\section{NES, MT-32 et puces PCM}
\capturedouble{nes}{mt32parts}{NES}{Roland MT-32 : parties}
Sur NES, les cartes des deux voix carrées dessinent leur rapport cyclique (12,5, 25, 50 ou 75~\%) et donnent le
volume sur 4 bits, l'enveloppe et le balayage de hauteur (« sweep ») ; le triangle et le bruit suivent, puis l'état
du canal DMC, qui joue les échantillons. Pour le MT-32, il n'y a pas de registres utiles à montrer : chaque partie
affiche le nom de son instrument, les touches jouées et l'état de ses « partiels », les briques de synthèse que le
module combine pour former un son.
\capturedouble{pcmchips-model1-portrait}{pcmchips-scsp-portrait}{YMW258 (Model 1/2)}{SCSP (Saturn, Model 2A/3)}
Les puces PCM des bornes Sega et de la Saturn jouent des dizaines d'échantillons à la fois. Chaque case montre une
voix active : son mode (PCM, FM ou bruit sur le SCSP), sa note, son panoramique et son niveau total (TL) en décibels.
